\section{Discussion}
\label{sec:discussion}

\para{Why are the biases so small?} We see three contributing factors. First, {\bf instruction tuning and scale weaken the demonstration prior.} At 1.5B, the base model is twice as sensitive to skew as its instruction-tuned sibling, both at 80\% (+8.9 vs. +4.9) and at 100\% (+29.5 vs. +15.6). The large effects reported for GPT-3~\citep{zhao2021calibrate} come from base models, which is consistent with this gap. Our design cannot separate model era from training recipe, however. Second, {\bf bias appears where the model is uncertain.} On \sst, where most models are above 90\% accurate, no manipulation moves predictions. On \mrpc, where accuracy is lowest and predictions are least stable, skew moves small models by 17--42 points. A simple account fits both observations: the demonstration label distribution acts as a prior, and this prior only wins when the evidence from the input is weak. As models become more competent and more confident (normalised entropy falls from 0.43--0.65 at $\le$1.5B to about 0.10 for Qwen at 7B and 14B), the prior has less room to act. Third, the four tasks are public and probably seen in pretraining, which may make models rely less on demonstrations than they would on a new task.

\para{Why anti-recency?} For models of 3B and above, two final $T$ demonstrations make $T$ \emph{less} likely, and primacy shows the mirror effect. We see two possible explanations. The model may expect labels to alternate, a gambler's-fallacy-like prior learned from roughly balanced data. Alternatively, this is a side effect of our design: with balanced counts, placing two $T$ demonstrations at the end leaves the middle positions slightly short of $T$ (only 38\% of the first two slots are $T$). With four final $T$ demonstrations the effect turns positive in small models but stays negative in Qwen2.5-7B. The recency bias described in the earlier literature therefore appears only when the run is long or the model is small. Separating the two explanations requires varying run length and alternation structure directly, which we leave to future work.

\para{Length is not a label cue.} Even with a 15--29 word gap between $T$ and non-$T$ demonstrations, no model learned to map length to label, either on average or on long test inputs. This fits the view of \citet{min2022rethinking} that models use demonstrations mostly for label space and format, rather than for fine-grained input--label statistics.

\para{Practical implications.} For instruction-tuned models of 3B and above on standard tasks, the specific labels of a balanced-looking demonstration set matter less than ordinary demonstration-to-demonstration variance. Practitioners will gain more from averaging over several demonstration sets, or from choosing good demonstrations, than from tuning label order or matching input lengths. Frequency skew still matters in three cases: small models, base models and ambiguous tasks. In these cases, balancing labels or applying contextual calibration~\citep{zhao2021calibrate,fei2023mitigating} remains worthwhile, and demonstration sets with a single label should be avoided.

\subsection{Limitations}
\label{sec:limitations}

\para{Model coverage.} We study one dense ladder (\qwen) plus two Llama models. The 7B and 8B models are 4-bit AWQ checkpoints, so quantisation is confounded with size at those points. Our largest model has 14B parameters; we do not know whether the trends continue to 70B and beyond.

\para{Scale statistics.} Slope CIs reflect item-sampling noise only. With five sizes, the evidence about variation across models is weak, and the Spearman correlations rest on $n=5$.

\para{Design choices.} We implemented label-position correlation as prompt-level primacy. \recency and \primacy slightly unbalance the labels in the middle positions, and in \agnews, \textsc{recency4} raises $T$'s share to 40\%. We used a single raw-completion template with $k=10$ and fixed label words; chat-formatted prompts or other values of $k$ may behave differently.

\para{Sample size and data.} With 200 test items per task, the CIs on the 4-task mean are about $\pm$1--2 points, so effects below about 2 points are hard to resolve. All four tasks are public benchmarks that are probably in pretraining data, which may deflate bias estimates relative to new tasks. The dose--response set covers only five models (not 3B, 14B or \llamaeight) because of limited GPU time.

\para{Broader impact.} This work measures robustness and does not introduce new capabilities. Its main risk is over-generalisation: our null results for modern models do not show that few-shot classifiers are safe from demonstration bias on new, ambiguous or high-stakes tasks, where we find the largest effects.
